% !TeX root = euclid-without-words.tex
% !TeX Program=pdfLaTeX
%
% Chord angles
%
\begin{minipage}[t]{.48\textwidth}
\hspace{4em}
\begin{tikzpicture}[scale=.8]
% Circle goes from (a) to (b)
\coordinate (a) at (0,0);
\coordinate (b) at (6,0);
% Line containing lower points is (c) -- (d)
\coordinate (c) at (0,-2);
\coordinate (d) at (6,-1.5);
% Line containing upper points is (e) -- (f)
\coordinate (e) at (0,1.5);
\coordinate (f) at (6,4);
% Name the upper and lower lines
\path [name path=chord1] (c) -- (d);
\path [name path=chord2] (e) -- (f);
% Draw a circle whose center is half-way between (a) and (b) through (a)
\coordinate (mid) at ($ (a)!.5!(b) $);
\node [draw,thick,circle through=(a),name path=circ] at (mid) {};
% Get intersections of the upper and lower lines with the circle
\path [name intersections={of=circ and chord1,by={i1,i2}}];
\path [name intersections={of=circ and chord2,by={i3,i4}}];
% Name the two chords
\path [name path=x1] (i1) -- (i3);
\path [name path=x2] (i2) -- (i4);
% Get their intersection
\path [name intersections={of=x1 and x2,by={p}}];
% Draw four lines from intersection to circle
\draw[thick,red] (i1) -- (p) -- 
  (i4) node[red,below,xshift=7pt,yshift=-7pt] {$\bm{\beta}$} -- cycle;
\draw[thick,blue] (i2) node[blue,above,xshift=-15pt,yshift=9pt] {$\bm{\gamma}$} -- (p) -- (i3) -- cycle;
\node[thick,violet,below,xshift=2pt,yshift=-4pt] at (p) {$\bm{\alpha}$};
% Display formula in color
\begin{scope}[xshift=1.8cm,yshift=-1.2cm]
\node[anchor=base,violet] at (0,0) {$\bm{\alpha}$};
\node[anchor=base] at (.5,0) {$\bm{=}$};
\node[anchor=base,red] at (1.0,0) {$\bm{\beta}$};
\node[anchor=base] at (1.5,0) {$\bm{+}$};
\node[anchor=base,blue] at (2,0) {$\bm{\gamma}$};
\end{scope}
\end{tikzpicture}
\end{minipage}
